% GENERATED FILE. Edit canonical lessons, then run npm run generate:books.
% canonical-source-hash: fnv1a64:1061a1d15a327f70
% canonical-lessons: ES-C439-anden, ES-C439-trasladar, ES-C439-cerrado, ES-C439-cortar, ES-C439-tirar, ES-R439-repaso-via, ES-C439-sintesis-via

\chapter{Andenes Y Cierres}
\label{ch:es-via}
\hlchaptermodality{439}

\begin{chapteropening}
\textbf{By the end of this chapter:} \emph{I can read the signs that decide where I may walk: a platform number, a shop shut for now, a street severed, and an office that has moved.}

Read the notice or message this chapter's words exist for, then say the same thing back.
\end{chapteropening}

\section[el andén]{el andén — where you wait for the train}

\label{lesson:ES-C439-anden}



\begin{quote}
Say \textbf{el tren}, \textbf{la estación} and \textbf{la parada}. You are inside the station and the train has not come. Name the strip of concrete you are standing on.
\end{quote}

\subsection*{You'll want to know: el andén}
\textbf{el andén} — the platform.

\emph{El tren sale del andén 4} — the train leaves from platform 4. Announcements put the number after the word, never before, and they say it fast: \emph{andén cuatro}.

The written accent tells you where to put the weight: \textbf{an-DÉN}, on the last syllable. Without it the word would be read \emph{AN-den}, and the stress is the part a listener uses to recognise it.

\begin{cousinweb}
It does not come apart, and this time that is not a trick. Nobody knows where \textbf{andén} comes from.

The tempting guess is \textbf{andar}, to walk, because a platform is a place you walk along. Dictionaries do not accept it: the endings do not line up, and the word turns up first in senses that have nothing to do with walking. Other proposals have been made and none has carried the field.

So the honest entry is \emph{origin uncertain}, and that is worth meeting on purpose. Any book that produces an etymology for every word is inventing some of them. Spanish has a few hundred words like this, and the right response is to learn the word and let the history stay open.
\end{cousinweb}

\subsection*{Your turn}
Say these aloud:

\begin{itemize}
\raggedright
  \item \emph{el andén}
  \item the train leaves from platform four
  \item the difference between \emph{la parada} and \emph{el andén}
\end{itemize}

\begin{tcolorbox}[breakable,colback=teal!4,colframe=teal!35!black,title={Before you move on}]
The platform? (\textbf{El andén}.) Which syllable is stressed? (\textbf{The last}.) Where does the word come from? (\textbf{Nobody knows for certain}.)
\end{tcolorbox}
\section[trasladar]{trasladar — carried across}

\label{lesson:ES-C439-trasladar}



\begin{quote}
Say \textbf{la oficina}, \textbf{el piso} and \textbf{el andén}. All three are places a thing or a person can be taken \textbf{from} one of and \textbf{to} another.
\end{quote}

\subsection*{You'll want to know: trasladar}
\textbf{trasladar} — to move something or someone from one place to another. \textbf{trasladarse} — to move yourself, to relocate.

\emph{Han trasladado la oficina al centro} — they have moved the office to the centre. \emph{Me traslado a Sevilla en enero} — I am relocating to Seville in January.

This is the formal word, the one on a letter from a bank or a clinic. For moving house among friends you would hear \emph{mudarse}, but the notice on the door will say \emph{trasladado}.

\begin{cousinweb}
\textbf{tras-} from Latin \emph{trans}, across, plus \textbf{-ladar} from \emph{lātus}, carried.

\noindent
\begin{tabularx}{\linewidth}{@{}>{\raggedright\arraybackslash}X>{\raggedright\arraybackslash}X@{}}
\toprule
\textbf{piece} & \textbf{meaning} \\
\midrule
tras- & across \\
-lad- & carried \\
\bottomrule
\end{tabularx}

\textbf{Carried across.} English has the same two pieces bolted together in \textbf{translate}, which is a text carried across into another language, and in \textbf{relate}, \textbf{collate} and \textbf{dilate}.
\end{cousinweb}

\begin{culture}
\emph{Lātus} is the odd piece, and it is worth a minute.

Latin's verb for carrying was \emph{ferre}. Its past participle was not \emph{ferred} or anything built from \emph{fer-}; it was \textbf{lātus}, from a completely different root that had been absorbed into the same verb long before anyone wrote it down. A verb whose forms come from two unrelated stems is called \textbf{suppletive}, and English has the famous one: \emph{go} and \emph{went}, where \emph{went} was borrowed whole from \emph{wend}.

So \textbf{trasladar} carries a fossil. The \emph{-lad-} in it is the past participle of a Latin verb whose present tense left no trace in the Spanish word at all.
\end{culture}

\subsection*{Your turn}
Say these aloud:

\begin{itemize}
\raggedright
  \item \emph{trasladar}
  \item they have moved the office
  \item I am relocating to a new flat in January
\end{itemize}

\begin{tcolorbox}[breakable,colback=teal!4,colframe=teal!35!black,title={Before you move on}]
To move something across? (\textbf{Trasladar}.) What do the two pieces mean? (\textbf{Across} and \textbf{carried}.) Which English word is built from the same two? (\textbf{Translate}.)
\end{tcolorbox}
\section[cerrado]{cerrado — the word on the glass}

\label{lesson:ES-C439-cerrado}



\begin{quote}
Say \textbf{cerrar} and \textbf{la puerta}. The office has been \textbf{trasladada} and you are standing in front of the old door. Name what the card hanging on it says.
\end{quote}

\subsection*{You'll want to know: cerrado}
\textbf{cerrado} — closed.

\emph{CERRADO POR VACACIONES} — closed for holidays. \emph{La piscina está cerrada los lunes} — the pool is closed on Mondays. It agrees like any adjective: \emph{cerrado}, \emph{cerrada}, \emph{cerrados}, \emph{cerradas}, matching the thing that is shut.

It also stretches past doors. \emph{Una curva cerrada} is a tight bend; \emph{un acento cerrado} is a thick accent you struggle to follow. The common thread is \textbf{not letting anything through}.

\begin{cousinweb}
\textbf{cerrar} plus \textbf{-ado}, the ending that turns an \emph{-ar} verb into its past participle. Nothing exotic: \emph{hablar} gives \emph{hablado}, \emph{cerrar} gives \emph{cerrado}.

Underneath sits Latin \textbf{claudere}, to shut, which you met under \emph{cerrar}. It is still visible in \textbf{la clausura}, \textbf{excluir} and English \textbf{closure}, \textbf{clause} and \textbf{claustrophobia} — a fear named for a shut room.
\end{cousinweb}

\begin{grammarlens}[title={Grammar Lens: two participles, two shapes}]
\textbf{abrir} gave you \textbf{abierto}, which does not follow the pattern. \textbf{cerrar} gives \textbf{cerrado}, which does. The pair sits on opposite sides of the same shop door, and their shapes do not match.

\noindent
\begin{tabularx}{\linewidth}{@{}>{\raggedright\arraybackslash}X>{\raggedright\arraybackslash}X@{}}
\toprule
\textbf{verb} & \textbf{participle} \\
\midrule
cerrar & cerrado — regular \\
abrir & \textbf{abierto} — irregular \\
\bottomrule
\end{tabularx}

The regular ones you build; the irregular ones you learn one at a time, and there are only about a dozen that matter. Spanish put the irregularity on the oldest and most-used verbs, which is the usual bargain a language makes: the words you say constantly are the ones that resist tidying.
\end{grammarlens}

\subsection*{Your turn}
Say these aloud:

\begin{itemize}
\raggedright
  \item \emph{cerrado}, then \emph{cerrada}
  \item the office is closed because they have moved it
  \item \emph{abierto} and \emph{cerrado}, and which one is the irregular shape
\end{itemize}

\begin{tcolorbox}[breakable,colback=teal!4,colframe=teal!35!black,title={Before you move on}]
Closed? (\textbf{Cerrado}.) Feminine? (\textbf{Cerrada}.) Is it a regular participle? (\textbf{Yes — \emph{-ar} gives \emph{-ado}}.)
\end{tcolorbox}
\section[cortar]{cortar — and what a cut road means}

\label{lesson:ES-C439-cortar}



\begin{quote}
Say \textbf{corto} and \textbf{el cuchillo}. One is a shape, one is a tool, and the same idea is inside both. Name the verb.
\end{quote}

\subsection*{You'll want to know: cortar}
\textbf{cortar} — to cut.

\emph{Corta el pan} — cut the bread. \emph{Me he cortado} — I have cut myself. \emph{Se ha cortado la llamada} — the call has dropped, which Spanish sees as a line cut rather than a call lost.

\begin{cousinweb}
\textbf{corto} and \textbf{cortar} are the same Latin word, one taken as a description and one as an action. \emph{Curtus} meant shortened, cut off.

\noindent
\begin{tabularx}{\linewidth}{@{}>{\raggedright\arraybackslash}X>{\raggedright\arraybackslash}X@{}}
\toprule
\textbf{word} & \textbf{the same idea as} \\
\midrule
corto & the result of the cut \\
cortar & the act of cutting \\
\bottomrule
\end{tabularx}

English keeps \emph{curtus} in \textbf{curt}, a reply cut short, and in \textbf{curtail}. The scissors in a barber's window and the adjective for a short skirt are relatives.
\end{cousinweb}

\begin{culture}
Now the meaning that catches travellers out. On a road or a rail line, \textbf{cortado} does not mean damaged. It means \textbf{closed}.

\begin{quote}
\textbf{VÍA CORTADA POR OBRAS}
\end{quote}

A \textbf{vía cortada} is a line out of service. A \textbf{calle cortada} is a street shut to traffic, often for a market or a race, with nothing wrong with it at all.

Set it beside \textbf{cerrado}, which you met a moment ago, because the two are not interchangeable. A shop is \textbf{cerrado}: someone shut it and will open it again. A road is \textbf{cortada}: the through-route has been severed, and you must go round. Spanish picks the verb by \emph{how} the way is blocked, and a sign that says \emph{cortada} is telling you there is no waiting for it to reopen today.
\end{culture}

\subsection*{Your turn}
Say these aloud:

\begin{itemize}
\raggedright
  \item \emph{cortar}
  \item the street is closed for works
  \item why a shop is \emph{cerrada} but a street is \emph{cortada}
\end{itemize}

\begin{tcolorbox}[breakable,colback=teal!4,colframe=teal!35!black,title={Before you move on}]
To cut? (\textbf{Cortar}.) What does \emph{VÍA CORTADA} mean? (\textbf{The line is closed}.) Which adjective you already had comes from the same Latin word? (\textbf{Corto}.)
\end{tcolorbox}
\section[tirar]{tirar — the door that says pull}

\label{lesson:ES-C439-tirar}



\begin{quote}
Say \textbf{la caja} and \textbf{la puerta}. The street was \textbf{cortada}, you carried an empty box round the long way, and now you want rid of it. Name what you do with the box.
\end{quote}

\subsection*{You'll want to know: tirar}
\textbf{tirar} — to throw, and above all \textbf{to throw away}.

\emph{Tira la caja} — throw the box out. \emph{No tires eso} — do not throw that away. \emph{Tirar la basura} is the everyday phrase for taking the rubbish out, and the bin itself is where things \textbf{se tiran}.

\begin{culture}
Now the sign, and it surprises every English speaker exactly once.

\begin{quote}
\textbf{TIRAR} on a door means \textbf{PULL}.
\end{quote}

The opposite side says \textbf{EMPUJAR}, push. A reader who knows \emph{tirar} as \emph{throw} reads the door as an instruction to hurl something, and gets it backwards.

The two senses are one idea seen from each end. \textbf{Tirar} is about applying force along a line away from where the thing sits: a box goes out of your hand, and a door comes out of its frame toward you. Spanish did not split the word; English happens to use two.

\noindent
\begin{tabularx}{\linewidth}{@{}>{\raggedright\arraybackslash}X>{\raggedright\arraybackslash}X@{}}
\toprule
\textbf{the sign} & \textbf{what to do} \\
\midrule
TIRAR & pull \\
EMPUJAR & push \\
\bottomrule
\end{tabularx}

Read the door before you lean on it. The mistake is harmless and universal, and knowing it in advance saves you making it in a doorway with people behind you.
\end{culture}

\begin{cousinweb}
It does not come apart, and like \textbf{el andén} at the start of this chapter its history is unsettled. A Germanic verb meaning to tear, reaching Spanish through the Goths, is the usual suggestion; an origin in the sound of a snapping cord has also been argued. Neither is agreed.

Two words in one chapter with no recoverable history is not unusual. The very common, very short, very old verbs are the ones whose trail goes cold soonest, because they were worn down long before anyone was writing Spanish at all.
\end{cousinweb}

\subsection*{Your turn}
Say these aloud:

\begin{itemize}
\raggedright
  \item \emph{tirar}
  \item throw the box away, the street is closed anyway
  \item what to do at a door marked \emph{TIRAR}
\end{itemize}

\begin{tcolorbox}[breakable,colback=teal!4,colframe=teal!35!black,title={Before you move on}]
To throw away? (\textbf{Tirar}.) What does \emph{TIRAR} mean on a door? (\textbf{Pull}.) And the sign on the other side? (\textbf{Empujar}, push.)
\end{tcolorbox}
\section[(andén, trasladar, cerrado, cortar …]{(andén, trasladar, cerrado, cortar, tirar) — from cold}

\label{lesson:ES-R439-repaso-via}



\begin{quote}
Close the lessons before this one. Say all five: \emph{the platform}, \emph{to move something across}, \emph{closed}, \emph{to cut}, \emph{to throw away}.
\end{quote}

\begin{grammarlens}[title={Grammar Lens: three signs that stop you, and they are not the same}]
English puts \emph{closed} on all three of these. Spanish does not, and the word it picks tells you what to do next.

\noindent
\begin{tabularx}{\linewidth}{@{}>{\raggedright\arraybackslash}X>{\raggedright\arraybackslash}X@{}}
\toprule
\textbf{the sign} & \textbf{what it tells you} \\
\midrule
\textbf{CERRADO} & shut for now; come back later \\
\textbf{CORTADA} & the route is severed; go round \\
\textbf{TRASLADADO} & it is open, somewhere else \\
\bottomrule
\end{tabularx}

A card reading \emph{TRASLADADO} is the one worth recognising fastest, because it is the only one of the three with an address underneath it.
\end{grammarlens}

\subsection*{Your turn}
Say these aloud:

\begin{itemize}
\raggedright
  \item the train leaves from platform four
  \item throw the box away
  \item what you do at a door marked \emph{TIRAR}, and at one marked \emph{EMPUJAR}
\end{itemize}

\begin{grammarlens}[title={What you've built}]
\textbf{You can now read the signs that decide where you are allowed to walk.} A station, a shut shop, a severed street and an office that has gone somewhere else are four different pieces of bad news, and you can tell them apart at a glance instead of guessing from context.

You also met the limit of etymology. \textbf{el andén} and \textbf{tirar} have no recoverable history, and the honest entry for both is \emph{origin uncertain}. Every other word in this chapter and the last came apart; these two did not, and saying otherwise would mean making something up.
\end{grammarlens}

\begin{tcolorbox}[breakable,colback=teal!4,colframe=teal!35!black,title={Before you move on}]
To cut? (\textbf{Cortar}.) To throw away? (\textbf{Tirar}.) Closed, on a shop door? (\textbf{Cerrado}.) And on a street? (\textbf{Cortada}.)
\end{tcolorbox}
\section[Practice]{(la oficina que se ha mudado) — read it, then do it yourself}

\label{lesson:ES-C439-sintesis-via}



\begin{quote}
Say \emph{el andén} and \emph{cerrado}. You are going to meet four signs in a row, and every one of them is going to send you somewhere else.
\end{quote}

\subsection*{Your turn}
Read all four, then answer.

\begin{quote}
\textbf{1 — ANDÉN 2. SALIDA A LA CALLE.}
\end{quote}

\begin{quote}
\textbf{2 — CALLE MAYOR CORTADA AL TRÁFICO.}
\end{quote}

\begin{quote}
\textbf{3 — CERRADO DE 14:00 A 16:00.}
\end{quote}

\begin{quote}
\textbf{4 — OFICINA TRASLADADA AL NÚMERO 40.}
\end{quote}

Say these aloud:

\begin{itemize}
\raggedright
  \item which sign means the place still exists
  \item which sign is about a time and not a place
  \item what sign 2 tells you to do about your route
\end{itemize}

\begin{culture}
Sign 2 repays a second look. \textbf{Cortada al tráfico} names \emph{what} the cut applies to, and that qualifier is doing real work: a street can be cut to cars and still open to feet. Spanish puts the limit of the cut inside the same phrase, where English would need a second sentence to say who is still allowed through.

Sign 4 is the one that saves your afternoon: \textbf{trasladada}, agreeing with \emph{oficina}, followed by an address. Compare it with a card that said only \emph{CERRADO} — that one would leave you waiting on a pavement for a door that is never opening again.

And when you finally reach number 40, the glass door will say \textbf{TIRAR}. Pull it.
\end{culture}

\begin{tcolorbox}[breakable,colback=teal!4,colframe=teal!35!black,title={Before you move on}]
Now your turn, out loud and then in writing. Five sentences: which platform you came out at, what the street was like, why the first door was no use, where the office went, and what you had to do to the last door.

Then check yourself: did you make \textbf{cortada} and \textbf{trasladada} agree with the words they describe?
\end{tcolorbox}
